\SourceLeaf{43}widmete. Das dort nicht in voller Allgemeinheit ausgedrückte Gesetz in übersichtlicherer Schreibweise zu bestimmen und mit strengem Beweis aus seiner eigentlichen Quelle herzuleiten, schien mir der Mühe wert.

Zunächst will ich kurz einige Vorbemerkungen zu einem neuen Zeichen machen, das bei der Darstellung allgemeiner Reihenglieder sehr nützlich ist. Bei Reihen, besonders bei den komplizierteren, geschieht es oft, dass die ungeraden Glieder einem anderen Gesetz zu folgen scheinen als die geraden, und dass die Reihe dann gewöhnlich in zwei voneinander getrennte Reihen zerlegt wird. Da dies mit einer Verlängerung des Ausdrucks verbunden ist und sozusagen die Einheitlichkeit des Ausdrucks aufhebt, durch die allein sich eine unbestimmte Reihe behandeln lässt, würde es die Verwendung von Reihen in enge Grenzen einschließen, wenn man nicht auf irgendeine Weise Abhilfe schüfe. Das allgemeinere Verfahren, das ich nun vorschlagen möchte, scheint sich von selbst aus dem Verfahren zu ergeben, mit dem der einfachste mögliche Unterschied zwischen den geraden und ungeraden Gliedern einer Reihe, nämlich dass die einen positiv und die anderen negativ sind, in das allgemeine Glied selbst aufgenommen wird. Das allgemeine Glied einer Reihe sei $m^{\text{tus}} = Q$. Soll es für gerades $m$ positiv und für ungerades $m$ negativ sein, so stelle ich ihm den Faktor $(-1)^m$ voran. Sollen die Vorzeichen in umgekehrter Ordnung aufeinanderfolgen, so stelle ich $(-1)^{m+1}$ voran\HistoricalNote{*)}{*)}{Der berühmte\ \textsc{Kaestner} war so freundlich, diesem Zeichen seine Zustimmung nicht zu versagen, als ich es ihm vorschlug. Auf seine Autorität gestützt, habe ich es bereits mehrfach in analytischen Untersuchungen verwendet.}. Ebenso sei der allgemeine Ausdruck für das Glied $m^{\text{ti}}$ bei geradem $m$ $= P$ und bei ungeradem $= Q$, sodass $P$ einem ganz anderen Gesetz folgt als $Q$. Ich behaupte, dass das Glied $m^{\text{tum}}$ allgemein $\frac{1-(-1)^m}{2}Q+\frac{1-(-1)^{m+1}}{2}P$ ist. Denn ist $m$ gerade, so ergibt $(-1)^m=1$ den Wert $\frac{1-(-1)^m}{2}=0$ und $(-1)^{m+1}=-1$ den Wert $\frac{1-(-1)^{m+1}}{2}=1$. Dadurch fällt der erste Teil weg und der andere wird wie verlangt $=P$. Ist dagegen $m$ ungerade, so gilt $(-1)^{m+1}=1$ und daher $\frac{1-(-1)^{m+1}}{2}=0$, zugleich $(-1)^m=-1$ und daraus $\frac{1-(-1)^m}{2}=1$. Der zweite Teil verschwindet und der erste bleibt wie vorausgesetzt $=Q$. Um nur ein Beispiel anzuführen, das den Nutzen der Schreibweise deutlich macht, nehme ich an, was unten bewiesen wird, dass sich eine beliebige Potenz des Kosinus eines gegebenen Winkels, z.\,B.\ $\cos\varphi^n$, stets in eine Reihe nach Kosinus von Vielfachen desselben Winkels entwickeln lässt. Dabei gilt das Gesetz, dass die Reihe bei geradem $n$ nach geraden Vielfachen fortschreitet und ihr Glied $r^{\text{tus}}$ unter Weglassung des Koeffizienten $\cos 2r\varphi$ ist. Bei ungeradem $n$ schreitet sie dagegen nach ungeraden Vielfachen fort und ihr Glied lautet dann $r^{\text{tus}} = \cos(2r+1)\varphi$. Der allgemeine Ausdruck lässt sich daher wie folgt angeben. In der Reihe, in die $\cos\varphi^n$%
